%==============================================================================
% فصل اول: مقدمه، کلیات، بیان مسأله
%==============================================================================
\ChapterStart
\chapter{مقدمه، کلیات و بیان مسأله}
\label{ch:introduction}
\section{مقدمه}
\label{sec:ch1-intro}

پرفشاری خون یکی از مهم‌ترین بیماری‌های مزمن و از چالش‌های پایدار نظام سلامت است. این بیماری ممکن است برای مدت طولانی بدون نشانه آشکار باقی بماند، اما در صورت تشخیص دیرهنگام یا کنترل نامطلوب، زمینه بروز پیامدهای جدی قلبی، مغزی، کلیوی و عروقی را فراهم می‌کند. ازاین‌رو، مدیریت پرفشاری خون به تجویز دارو یا پیگیری‌های دوره‌ای محدود نمی‌شود و مستلزم مشارکت آگاهانه و مستمر بیمار در فرایند مراقبت است. توانایی بیمار برای پیگیری وضعیت خود و اتخاذ تصمیم‌های مناسب در زندگی روزمره، بخش مهمی از کنترل بیماری و پیشگیری از عوارض آن را تشکیل می‌دهد \cite{han2014hscp,ghaneigheshlagh2019persian,lee2022chronicmeta,aghajanloo2021selfcare}.
\par
خودمراقبتی در پرفشاری خون مجموعه‌ای از رفتارهای هدفمند و پیوسته است که پایش صحیح فشار خون در منزل، مصرف منظم داروهای تجویزشده، رعایت الگوی غذایی سالم، کاهش مصرف نمک و چربی، حفظ وزن مناسب، انجام فعالیت بدنی متناسب، مدیریت تنش‌های روانی و مراجعه به‌موقع برای دریافت مراقبت را دربر می‌گیرد. انجام مؤثر این رفتارها به آگاهی از ماهیت بیماری، نگرش مناسب، مهارت عملی و انگیزه کافی نیاز دارد. بااین‌حال، بی‌نشانه بودن بیماری، طولانی بودن دوره درمان، دشواری تغییر عادت‌های تثبیت‌شده و فراموشی توصیه‌ها می‌تواند استمرار خودمراقبتی را با مشکل روبه‌رو سازد. بنابراین، آموزش بیمار زمانی ارزشمند است که افزون بر انتقال اطلاعات، فهم، مهارت، مشارکت و تداوم رفتارهای مراقبتی را تقویت کند \cite{wilkinson2009selfcare,slemon2021selfcare,esmaili2020selfcare,han2014hscp,ghaneigheshlagh2019persian,rakhshani2024precede,abedi2020precede,ghorbani2021vark,nasirian2019group,suprayitno2023telenursing,kurt2022selfmgmt,desouza2016flashcard}.
\par
آموزش کارگاهی یکی از روش‌های رایج آموزش بیماران است که امکان ارتباط مستقیم میان آموزش‌دهنده و یادگیرنده را فراهم می‌آورد. در این روش، مطالب را می‌توان با سخنرانی کوتاه، پرسش و پاسخ، بحث گروهی و تمرین عملی ارائه کرد و ابهام‌های بیماران را در همان جلسه برطرف ساخت. تمرین اندازه‌گیری فشار خون، گفت‌وگو درباره موانع رعایت رژیم غذایی و بررسی راهکارهای پایبندی به درمان از جمله فعالیت‌هایی است که در محیط کارگاه قابلیت اجرا دارد. تعامل چهره‌به‌چهره و دریافت بازخورد فوری از نقاط قوت این شیوه است؛ بااین‌حال، محدود بودن زمان جلسات، وابستگی آموزش به حضور در مکان مشخص و کاهش احتمال مرور منظم مطالب پس از پایان کارگاه می‌تواند ماندگاری و کاربرد روزانه آموخته‌ها را محدود کند \cite{poulsen2015workshop,schwill2023workshop,slatyer2017mscr,yu2022empowerment,zhang2024psychotherapists,lee2022chronicmeta}.
\par
گسترش ابزارهای ارتباطی و دسترسی روزافزون به تلفن همراه، امکان تکمیل آموزش حضوری با روش‌های انعطاف‌پذیرتر را فراهم کرده است. خرده‌یادگیری یکی از این روش‌هاست که محتوای آموزشی را در واحدهای کوتاه، متمرکز و قابل‌فهم ارائه می‌کند. کوتاه بودن هر واحد آموزشی می‌تواند دریافت محتوا را در زمان‌ها و موقعیت‌های مختلف آسان‌تر سازد و مرور چندباره نکات کلیدی را امکان‌پذیر کند. در آموزش بیماران مبتلا به پرفشاری خون، پیام‌ها یا ویدئوهای کوتاه می‌توانند برای یادآوری شیوه صحیح اندازه‌گیری فشار خون، انتخاب غذای مناسب، فعالیت بدنی، مدیریت تنش و مصرف منظم داروها به‌کار روند. بااین‌حال، خرده‌یادگیری جایگزین کامل تعامل انسانی و آموزش عملی نیست و ارزش آن باید در پیوند با نیازهای بیمار، کیفیت محتوا و شیوه ارائه بررسی شود \cite{wang2020microlearning,monib2024microlearning,poorcheraghi2025microlearning,susman2024micropractice}.
\par
تلفیق آموزش کارگاهی با خرده‌یادگیری می‌تواند ظرفیت‌های مکمل این دو روش را در کنار یکدیگر قرار دهد. کارگاه، فرصت توضیح عمیق‌تر، تمرین مهارت و تعامل مستقیم را فراهم می‌سازد و محتوای کوتاه دیجیتال می‌تواند فاصله میان جلسات را با مرور و یادآوری هدفمند پوشش دهد. چنین ترکیبی از نظر آموزشی ممکن است به حفظ پیوستگی یادگیری و انتقال بهتر آموخته‌ها به موقعیت‌های واقعی زندگی کمک کند. بااین‌همه، جذابیت یا سهولت دسترسی به محتوای کوتاه به‌تنهایی برای اثبات بهبود عملکرد خودمراقبتی کافی نیست. اثربخشی این رویکرد باید در مقایسه با آموزش کارگاهی و با سنجش پیامدهای مرتبط با رفتار بیمار و کنترل فشار خون ارزیابی شود \cite{ghiyasvandian2023blended,rahbar2024socialmedia,munro2018blended,jennatfereidooni2026orem,rooddehghan2024ms,wang2020microlearning,schwill2023workshop,nourse2025heartfailure,adnan2022wellbeing,adam2023digital}.
\par
انتخاب این موضوع از آن جهت اهمیت دارد که بیماران مبتلا به پرفشاری خون برای مدیریت روزانه بیماری خود به آموزشی عملی، قابل‌استمرار و متناسب با محدودیت‌های زندگی روزمره نیاز دارند. در عین حال، کارکنان سلامت برای انتخاب شیوه آموزشی مناسب به شواهدی نیازمندند که ارزش افزوده روش‌های جدید را فراتر از آموزش حضوری روشن سازد. مشخص شدن اینکه افزودن خرده‌یادگیری به آموزش کارگاهی تا چه اندازه می‌تواند عملکرد خودمراقبتی را ارتقا دهد، به تصمیم‌گیری آگاهانه‌تر درباره طراحی برنامه‌های آموزش بیمار و تخصیص منابع آموزشی کمک می‌کند \cite{wang2020microlearning,schwill2023workshop,nourse2025heartfailure,rakhshani2024precede,abedi2020precede,ghorbani2021vark,nasirian2019group,suprayitno2023telenursing,kurt2022selfmgmt,desouza2016flashcard,lee2022chronicmeta}.
\par
بر همین اساس، پژوهش حاضر با هدف مقایسه اثربخشی آموزش خودمراقبتی به روش کارگاهی و آموزش کارگاهی تلفیق‌شده با خرده‌یادگیری بر عملکرد خودمراقبتی افراد مبتلا به پرفشاری خون انجام شد. در رویکرد تلفیقی، آموزش حضوری با ویدئوهای کوتاه و هدفمند همراه شد تا امکان مرور تدریجی مطالب و یادآوری رفتارهای اصلی خودمراقبتی فراهم شود. تمرکز این پژوهش بر مقایسه دو شیوه آموزشی، زمینه‌ای برای ارزیابی ارزش افزوده خرده‌یادگیری در کنار کارگاه و بررسی کاربرد آن در آموزش بیماران مبتلا به یک بیماری مزمن فراهم می‌آورد \cite{poulsen2015workshop,schwill2023workshop,zhang2024psychotherapists,slatyer2017mscr,ghiyasvandian2023blended,rahbar2024socialmedia,munro2018blended,jennatfereidooni2026orem,wang2020microlearning,monib2024microlearning,lee2022chronicmeta,rakhshani2024precede,abedi2020precede,ghorbani2021vark,nasirian2019group,ghaneigheshlagh2019persian,han2014hscp}.


\section{بیان مسأله}
\label{sec:problem-statement}
پرفشاری خون یکی از شایع‌ترین بیماری‌های مزمن و از مهم‌ترین عوامل قابل اصلاح در بروز بیماری‌های قلبی‌عروقی، سکته مغزی، نارسایی کلیوی و مرگ زودرس است \cite{han2014hscp,aghajanloo2021selfcare}. ماهیت اغلب بی‌علامت این بیماری سبب می‌شود بسیاری از مبتلایان تا زمان بروز عوارض، خطر موجود را به‌درستی درک نکنند یا ضرورت پیگیری مستمر درمان را جدی نگیرند \cite{rakhshani2024precede,abedi2020precede}. از سوی دیگر، کنترل پرفشاری خون تنها به تشخیص پزشکی و تجویز دارو محدود نیست، بلکه مستلزم مشارکت فعال و مداوم بیمار در مراقبت روزانه است \cite{han2014hscp,wilkinson2009selfcare}. ازاین‌رو، حتی در صورت دسترسی به خدمات درمانی، ضعف در رفتارهای خودمراقبتی می‌تواند دستیابی به کنترل مطلوب فشار خون و پیشگیری از عوارض را با مشکل روبه‌رو سازد.
\par
خودمراقبتی در پرفشاری خون مجموعه‌ای از تصمیم‌ها و رفتارهای روزمره را دربر می‌گیرد که بیمار برای مدیریت بیماری خود انجام می‌دهد. مصرف منظم و صحیح داروها، رعایت الگوی غذایی مناسب و محدودکردن مصرف نمک، فعالیت بدنی متناسب، کنترل وزن، پرهیز از دخانیات، مدیریت تنش‌های روانی، پایش فشار خون در منزل و مراجعه به‌موقع برای پیگیری درمان از مهم‌ترین این رفتارها هستند \cite{han2014hscp,wilkinson2009selfcare,slemon2021selfcare}. انجام پایدار چنین اقداماتی به آگاهی صرف وابسته نیست؛ بیمار باید ضرورت هر رفتار را درک کند، مهارت اجرای آن را به دست آورد، موانع شخصی و محیطی را بشناسد و در گذر زمان انگیزه خود را حفظ کند \cite{rakhshani2024precede,abedi2020precede,ghorbani2021vark}. بنابراین، آموزش بیمار زمانی می‌تواند به بهبود عملکرد خودمراقبتی بینجامد که افزون بر انتقال اطلاعات، فرصت تمرین، بازخورد، تکرار و یادآوری را نیز فراهم کند \cite{schwill2023workshop,susman2024micropractice}.
\par
آموزش کارگاهی یکی از روش‌های متداول برای آموزش خودمراقبتی به بیماران مزمن است. حضور هم‌زمان آموزش‌دهنده و بیماران، امکان پرسش و پاسخ، تبادل تجربه، رفع برداشت‌های نادرست و تمرین عملی مهارت‌ها را فراهم می‌سازد \cite{poulsen2015workshop,schwill2023workshop,zhang2024psychotherapists}. این ویژگی‌ها می‌توانند یادگیری را از دریافت منفعل اطلاعات فراتر ببرند و به شکل‌گیری مهارت و اعتماد بیمار کمک کنند \cite{slatyer2017mscr,zhang2024psychotherapists}. بااین‌حال، اثر آموزش کارگاهی ممکن است پس از پایان جلسات کاهش یابد. محدودبودن زمان جلسات، تراکم مطالب، تفاوت سرعت یادگیری بیماران، فراموشی تدریجی محتوا و نبود دسترسی مستمر به آموزش‌دهنده از عواملی هستند که انتقال آموخته‌ها به رفتار روزمره را دشوار می‌کنند \cite{poulsen2015workshop,schwill2023workshop,slatyer2017mscr}. همچنین، تکرار جلسات حضوری برای بیمار و نظام ارائه‌دهنده خدمت با محدودیت‌های زمانی، مکانی و اجرایی همراه است \cite{munro2018blended,rahbar2024socialmedia}.
\par
خرده‌یادگیری رویکردی است که محتوا را در واحدهای کوتاه، متمرکز و قابل استفاده در زمان‌های مختلف ارائه می‌کند \cite{wang2020microlearning,monib2024microlearning}. دسترسی به محتوای کوتاه از طریق تلفن همراه می‌تواند امکان مرور چندباره، دریافت تدریجی پیام‌ها و پیوند آموزش با موقعیت‌های واقعی زندگی را فراهم سازد \cite{monib2024microlearning,poorcheraghi2025microlearning,susman2024micropractice}. در بیماران مبتلا به پرفشاری خون، محتوای ویدئویی کوتاه می‌تواند به یادآوری رفتارهای مشخصی مانند شیوه صحیح اندازه‌گیری فشار خون، مصرف دارو، انتخاب غذای مناسب و فعالیت بدنی کمک کند \cite{poorcheraghi2025microlearning,desouza2016flashcard,han2014hscp}. بااین‌وجود، کوتاه‌بودن محتوا یا دیجیتال‌بودن شیوه ارائه به‌تنهایی تضمین‌کننده تغییر رفتار نیست. میزان ارتباط محتوا با نیاز بیمار، تداوم استفاده، توانایی کار با تلفن همراه، دسترسی به اینترنت و همراه‌شدن آموزش با تمرین و بازخورد می‌توانند بر اثربخشی این روش اثر بگذارند \cite{wang2020microlearning,susman2024micropractice,monib2024microlearning}.
\par
بررسی شواهد موجود نشان می‌دهد آموزش‌های حضوری و کارگاهی عموما می‌توانند دانش، خودکارآمدی و برخی رفتارهای خودمراقبتی را بهبود دهند \cite{poulsen2015workshop,schwill2023workshop,zhang2024psychotherapists,slatyer2017mscr}. روش‌های دیجیتال و خرده‌یادگیری نیز در افزایش آگاهی، یادآوری مطالب و تقویت پیامدهای نزدیک به یادگیری امیدوارکننده بوده‌اند \cite{wang2020microlearning,ghiyasvandian2023blended,rahbar2024socialmedia,munro2018blended}. بااین‌حال، بهبود دانش همواره به تغییر پایدار رفتار یا بهبود شاخص‌های بالینی منجر نمی‌شود \cite{wang2020microlearning,lee2022chronicmeta}. ناهمگونی در نوع بیماری، ویژگی‌های شرکت‌کنندگان، محتوای آموزشی، تعداد و طول جلسات، شیوه ارائه پیام‌های کوتاه، مدت پیگیری و ابزارهای سنجش نیز مقایسه نتایج مطالعات را دشوار کرده است \cite{lee2022chronicmeta,wang2020microlearning}. در بسیاری از پژوهش‌ها، روش ترکیبی با مراقبت معمول مقایسه شده یا آموزش دیجیتال جایگزین آموزش حضوری بوده است؛ بنابراین، سهم افزوده خرده‌یادگیری هنگامی که محتوای اصلی و آموزش کارگاهی در هر دو گروه یکسان باشد، هنوز به‌روشنی مشخص نیست \cite{ghiyasvandian2023blended,rahbar2024socialmedia,munro2018blended,nourse2025heartfailure}.
\par
در حوزه پرفشاری خون نیز مداخلات آموزشی متنوعی با استفاده از آموزش گروهی، الگوهای آموزش سلامت و ابزارهای دیجیتال بررسی شده‌اند و نتایج آنها در مجموع بر اهمیت آموزش بیمار دلالت دارد \cite{rakhshani2024precede,abedi2020precede,ghorbani2021vark,nasirian2019group,ramazankhani2019hbm,suprayitno2023telenursing,kurt2022selfmgmt,desouza2016flashcard}. بااین‌حال، شواهد کافی برای پاسخ به این پرسش وجود ندارد که آیا افزودن خرده‌یادگیری ویدئویی به یک برنامه کارگاهی مشخص، در مقایسه با همان برنامه کارگاهی به‌تنهایی، به بهبود بیشتر عملکرد خودمراقبتی و کنترل فشار خون منجر می‌شود یا خیر \cite{nourse2025heartfailure,wang2020microlearning,lee2022chronicmeta}. کمبود مقایسه مستقیم این دو روش، به‌ویژه در بیماران ایرانی و در شرایط واقعی ارائه خدمات، تصمیم‌گیری درباره ارزش آموزشی و اجرایی روش ترکیبی را دشوار می‌سازد \cite{rakhshani2024precede,abedi2020precede,ghorbani2021vark,nasirian2019group}. بدون چنین مقایسه‌ای نمی‌توان مشخص کرد که مزایای احتمالی روش ترکیبی ناشی از خرده‌یادگیری است یا صرفا از دریافت آموزش بیشتر و تماس مکرر با محتوای آموزشی حاصل می‌شود \cite{lee2022chronicmeta,ghiyasvandian2023blended,rahbar2024socialmedia}.
\par
مسئله اصلی پژوهش حاضر، روشن‌کردن همین اثر افزوده است. در این پژوهش، آموزش خودمراقبتی به روش کارگاهی با آموزش کارگاهی تلفیق‌شده با خرده‌یادگیری در افراد مبتلا به پرفشاری خون مقایسه شد. برنامه کارگاهی شامل چهار جلسه حضوری هفتگی شصت‌دقیقه‌ای بود و در روش تلفیقی، افزون بر همین جلسات، هشت ویدئوی آموزشی سه‌دقیقه‌ای طی چهار هفته در اختیار بیماران قرار گرفت. سنجش عملکرد خودمراقبتی و فشار خون پیش از مداخله و یک ماه پس از پایان آن، امکان بررسی تغییرات درون هر گروه و مقایسه میزان تغییر میان دو گروه را فراهم کرد \cite{han2014hscp,ghaneigheshlagh2019persian}. یکسان‌بودن بخش کارگاهی در دو گروه نیز زمینه را برای ارزیابی مشخص‌تر نقش خرده‌یادگیری به‌عنوان مکمل آموزش حضوری فراهم ساخت \cite{wang2020microlearning,munro2018blended}.
\par
بر این اساس، پرسش محوری آن بود که آیا تلفیق خرده‌یادگیری ویدئویی با آموزش کارگاهی، در مقایسه با آموزش کارگاهی به‌تنهایی، موجب بهبود بیشتر عملکرد خودمراقبتی و کاهش بیشتر فشار خون در بیماران مبتلا به پرفشاری خون می‌شود. فرض پژوهش بر این مبنا شکل گرفت که تکرار فاصله‌دار و دسترسی آسان به محتوای کوتاه می‌تواند ماندگاری آموخته‌ها و تبدیل آنها به رفتار روزمره را تقویت کند \cite{susman2024micropractice,jennatfereidooni2026orem,rooddehghan2024ms,wang2020microlearning} و در نتیجه، تغییرات مطلوب‌تری در پیامدهای خودمراقبتی و فشار خون پدید آورد \cite{yu2022empowerment,rakhshani2024precede,poorcheraghi2025microlearning}. پاسخ به این پرسش می‌تواند برای انتخاب شیوه آموزشی متناسب با نیاز بیماران، طراحی برنامه‌های قابل اجرا در مراکز درمانی و استفاده هدفمند از فناوری‌های در دسترس در مدیریت بیماری‌های مزمن سودمند باشد.

\section{اهمیت و ضرورت پژوهش}
\label{sec:significance}
پرفشاری خون به عنوان یک بیماری مزمن و شایع، یکی از مهم‌ترین چالش‌های نظام‌های سلامت در سراسر جهان است. این بیماری که به عنوان "قاتل خاموش" شناخته می‌شود، به دلیل عدم وجود علائم اولیه اغلب دیر تشخیص داده می‌شود و می‌تواند منجر به عوارض جدی همچون سکته قلبی و مغزی، نارسایی کلیوی و مشکلات عروقی گردد. بر اساس گزارش سازمان جهانی بهداشت، پرفشاری خون عامل اصلی بیش از ۷.۵ میلیون مرگ در جهان است.
\par
مفهوم "خود مراقبتی" به عنوان ابزار کلیدی در توانمندسازی بیماران مطرح است. رفتارهایی مانند پایش فشار خون، مصرف صحیح داروها، رعایت رژیم غذایی مناسب (مانند \lr{DASH})، انجام فعالیت‌های بدنی و مدیریت استرس نقش اساسی در کنترل بیماری دارد. آموزش به روش کارگاهی به دلیل تعامل مستقیم چهره‌به‌چهره بسیار موثر است اما محدودیت‌های زمانی و مکانی دارد. آموزش‌های مجازی در قالب "خرده یادگیری" (ارائه محتوا در قالب ویدیوهای کوتاه حدود ۳ دقیقه‌ای) می‌تواند پویایی و تداوم آموزش را افزایش داده و به عنوان مکمل قدرتمند به کار گرفته شود. این مطالعه اثربخشی روش کارگاهی سنتی را در مقایسه با روش ترکیبی (کارگاهی \+ خرده یادگیری ویدئویی) بررسی خواهد کرد.


\section{اهداف پژوهش}
\label{sec:objectives}

\subsection{هدف کلی}
مقایسه اثر‌بخشی آموزش خودمراقبتی به روش کارگاهی و روش کارگاهی تلفیق‌شده با
خرده‌یادگیری بر عملکرد خودمراقبتی افراد مبتلا به پرفشاری خون.

\subsection{اهداف اختصاصی}
\begin{enumerate}[label=\arabic*.]
  \item تعیین و مقایسه میانگین نمره عملکرد خودمراقبتی قبل و بعد از مداخله در
        هر یک از گروه‌های مطالعه.
  \item مقایسه میانگین نمره عملکرد خودمراقبتی پس از مداخله بین دو گروه.
  \item تعیین و مقایسه میانگین فشار خون سیستولیک و دیاستولیک قبل و بعد از
        مداخله در هر یک از دو گروه و مقایسه تغییرات بین دو گروه.
\end{enumerate}

\subsection{اهداف کاربردی}
ارائه الگوی آموزشی عملی برای آموزش بیماران مبتلا به بیماری‌های مزمن، به‌ویژه
پرفشاری خون، به‌منظور کمک به تیم درمان در انتخاب روش‌های آموزشی مؤثرتر.

\section{پرسش‌های پژوهش}
\label{sec:research-questions}
\begin{enumerate}[label=\arabic*.]
  \item آیا آموزش خودمراقبتی به روش کارگاهی بر عملکرد خودمراقبتی بیماران مبتلا
        به پرفشاری خون تأثیر دارد؟
  \item آیا آموزش خودمراقبتی به روش کارگاهی تلفیق‌شده با خرده‌یادگیری بر عملکرد
        خودمراقبتی بیماران مبتلا به پرفشاری خون تأثیر دارد؟
  \item آیا بین دو روش آموزشی از نظر اثر بر عملکرد خودمراقبتی تفاوت معنادار
        وجود دارد؟
\end{enumerate}

\section{فرضیات پژوهش}
\label{sec:hypotheses}
\begin{enumerate}[label=\arabic*.]
  \item آموزش خودمراقبتی به روش کارگاهی موجب بهبود معنادار عملکرد خودمراقبتی
        بیماران می‌شود.
  \item آموزش خودمراقبتی به روش کارگاهی تلفیق‌شده با خرده‌یادگیری موجب بهبود
        معنادار عملکرد خودمراقبتی بیماران می‌شود.
  \item اثر‌بخشی روش کارگاهی تلفیق‌شده با خرده‌یادگیری بر عملکرد خودمراقبتی
        بیشتر از روش کارگاهی تنها است.
\end{enumerate}

